% ===============================================================
\section{Summary of the program modifications (Task 2.6)}
\label{sec:codechanges}
% ===============================================================

This section lists, program by program, the modifications made to the previous version of the code, and
all the changes made to the basic structure of the example code. The complete programs are in
Appendix~\ref{app:code}. Every modified line is also marked in the code by a comment.

\subsection{Changes to the basic structure of the example code (all GA programs)}
\begin{enumerate}[label=\textbf{S\arabic*.}]
  \item \textbf{Bug fix -- \code{import random}.} In \code{CodeP2.2F26}, \code{import random} replaces the
        function \code{random()} imported in \code{CodeP2.1F26} by the module, so the combined program
        stops at the initialisation with a \code{TypeError}. It was replaced by
        \code{from random import random, seed}.
  \item \textbf{Data corrections.} Two values of \code{CodeP2.1F26} were corrected to match the Appendix
        (Table~\ref{tab:datacheck}).
  \item \textbf{Evolution loop placed in a function.} The skeleton code from ``INITIALIZING PARAMETERS'' to
        the final print was placed in \code{run\_GA(n1i,\dots,n5i, NGEN, MUT, MFRAC, rseed, target, verbose)}.
        This allows many runs (initial guesses, mutation amplitudes, seeds, NGEN) in the same notebook.
        Inside the function, the statements of the loop (fitness, selection, mating, mutation, averaging,
        minimum search, time averages) are those of the skeleton, line for line.
  \item \textbf{Parameters instead of literals.} The mutation amplitude (\code{0.085} in the code) is the
        argument \code{MUT}, and the selection fraction \code{MFRAC} ($=0.5$ in the skeleton) is an argument.
  \item \textbf{Reproducibility.} The skeleton seeds Python's \code{random} once but never seeds
        \code{numpy.random}, which is used for selection, mating and mutation. Every run therefore gave
        different results. \code{run\_GA} now calls \code{seed(rseed)} and \code{numpy.random.seed(rseed)},
        so any run of this report can be reproduced exactly from its seed.
  \item \textbf{Diagnostics returned.} The function returns the histories (\code{aFerrmeanavgn},
        \code{n1avg}, \dots), the generation of the minimum (\code{kmin}), the first generation below the
        target (\code{kconv}) and the number of survivors of each generation (\code{nkeephist}). None of
        these quantities affects the algorithm.
  \item \textbf{Figures.} The skeleton plots are kept (\code{plt.loglog()}, same legends). The
        \code{yticks} were extended down to $10^{-4}$ so that $n_1\approx10^{-3}$ is visible. Additional
        report-quality figures are saved as PDF files.
\end{enumerate}

\subsection{Program-by-program modifications}

\paragraph{\code{P2\_Task2.1} (from \code{CodeP2.1F26}).} S2; a print of the number of rows and of one
element; a cell that checks the data against the Appendix; a cell for the log-log plot and power-law
fits of Task~2.1(b).

\paragraph{\code{P2\_Task2.2} (from \code{CodeP2.1F26} + \code{CodeP2.2F26}).} S1--S7. The error
function of the skeleton (Eq.~(1.3)) is unchanged. Added cells: first run with the skeleton guesses,
deterministic reference, initial-guess study, best-fit plots and statistics, landscape.

\paragraph{\code{P2\_Task2.3} (copy of \code{P2\_Task2.2}).} Modifications (i)--(vii) of the statement:
\begin{enumerate}[label=(\roman*)]
  \item the red-text data lines (32 rows with $P=9.5$ and $7.0$~kPa, $\gamma=0$ and $1.71$) are
        uncommented, giving 77 rows;
  \item \code{ND = 77}, \code{NS = 77};
  \item the error function becomes Eq.~(1.7) in both locations (fitness loop and population-average loop),
        and the prediction used in the plots becomes Eq.~(1.5):
\begin{lstlisting}
gen_en = 9.8   # earth-normal gravitational acceleration, m/s2
Ferr[i] = n[i][0]*lydata[i][0] + math.log(n[i][1]) + n[i][2]*lydata[i][1]
Ferr[i] = Ferr[i] + n[i][3]*math.log( ydata[i][2] + n[i][4]*gen_en*ydata[i][3] ) + n[i][5]*lydata[i][4]
...
Ferravgn[i] = -1.*lydata[i][0] + math.log(n1avg[k]) + n2avg[k]*lydata[i][1]
Ferravgn[i] = Ferravgn[i] + n3avg[k]*math.log( ydata[i][2] + n4avg[k]*gen_en*ydata[i][3] ) + n5avg[k]*lydata[i][4]
...
qpred = n1*row[1]**n2 * (row[2] + n4*gen_en*row[3])**n3 * row[4]**n5
\end{lstlisting}
  \item the mating/mutation lines of $n_4$ and $n_5$ are uncommented (with amplitude \code{MUT});
  \item the final print-out lists $n_1$--$n_5$ (\code{ENDING}, \code{MINUMUM}, \code{TIME AVG});
  \item the legend of the first plot is replaced by the five-constant legend and $n_4$, $n_5$ are plotted;
  \item the default initial guesses are $0.000476,\ 3.028,\ 0.2249,\ 1.054,\ 0.217$.
\end{enumerate}
Added: the 15 validation points of the Appendix (\code{vdata}, used only to test the fit), the study
over initial guesses and the multi-stage refinement, the surface plot, and the profiles of
Fig.~\ref{fig:t23prof}.

\paragraph{\code{P2\_ET1} (new).} Reads the best-fit constants of Task~2.3 and evaluates the fractional
change ratios with finite steps, following steps 1--3 of the statement. It also computes the analytical
logarithmic derivatives and the limiting input inaccuracies.

\paragraph{\code{P2\_ET2} (copy of \code{P2\_Task2.3}).} The data rows are reorganised as
$[T_w-T_{sat},\ \qpp,\ g,\ \gamma,\ P]$ and \code{lydata} is recomputed. Because column~0 is now
$\ln\dTsat$ and column~1 is $\ln\qpp$, the error statements of Task~2.3 compute exactly the error of
Eq.~(ET2.1). Only the initial guesses (negative $\nh3$, $\nh5$: the multiplicative mutation keeps the
sign), the labels and the analysis cells were changed.

\paragraph{\code{P2\_Task2.4-2.5} (copy of \code{P2\_Task2.2}).}
\begin{itemize}
  \item The data are those of \code{CodeP2.3F26} (122 points $[T_{a,in},T_{a,out},\dot q]$). A seeded random
        permutation splits them into a training set (\code{ydata}, $N_D=98$, $80.3\,\%$) and a validation set
        (\code{vdata}, $N_V=24$, $19.7\,\%$). The GA sees only the training set, and $N_S=N_D=98$.
  \item The error function was changed in the two locations, Eqs.~(2.4)--(2.5), with $n_1=(UA)_e$,
        $n_2=(UA)_c$, $n_3=\alpha$. The model Eq.~(2.6) is coded once and used in both locations and in the plots:
\begin{lstlisting}
sigma = 5.67e-8     # Stefan-Boltzmann constant, W m-2 K-4
hcref = 80.0        # W m-2 K-1
def hrad(Tain, Taout):
    Ti = Tain + 273.15
    To = Taout + 273.15
    return sigma*(Ti + To)*(Ti**2 + To**2)
def qdot_model(n1, n2, n3, Tain, Taout):
    UAcr = n2*(1. + n3*hrad(Tain, Taout)/hcref)
    return UAcr*( n1*(Tain - Taout)/(n1 + UAcr) )
...
Ferr[i] = ydata[i][2] - qdot_model(n[i][1], n[i][2], n[i][3], ydata[i][0], ydata[i][1])
aFerr[i] = abs(Ferr[i])/abs(ydata[i][2])  #- absolute fractional error
...
Ferravgn[i] = ydata[i][2] - qdot_model(n1avg[k], n2avg[k], n3avg[k], ydata[i][0], ydata[i][1])
aFerravgn[i] = abs(Ferravgn[i])/abs(ydata[i][2])
\end{lstlisting}
  \item The mean absolute error of the validation set is computed at each generation with the
        population-average constants (\code{aFerrval}). It is used for monitoring only, not for selection.
  \item The plotting section was rewritten for the log-log plots of predicted versus data $\dot q$ (training
        and validation) and for the surface plot of $\dot q(T_{a,in},T_{a,out})$.
\end{itemize}
